\section{Arbitrary regular-sequence length and characteristic-zero fields}
\label{sec:characteristic-zero}

The new input of this paper is the Artinian theorem over \(\C\). We now
apply the published Artinian reduction and record the coefficient-field
argument needed for Corollary~\ref{cor:characteristic-zero}.

\begin{proof}[Proof of Corollary~\ref{cor:characteristic-zero}]
If \(I=S_F\), the prescription gives \(J=S_F\), and all quotient Betti
numbers vanish. We may therefore assume that \(I\) is proper.

First let \(F=\C\). After a linear change of coordinates, one may arrange
that
\[
 f_1,\ldots,f_c,x_{c+1},\ldots,x_n
\]
is a regular sequence; the list of variables is empty when \(c=n\).
Write \(\overline f_i\) for the images of the \(f_i\) in
\(\overline S=\C[x_1,\ldots,x_c]\). They form a full regular sequence
of degrees \(a_1,\ldots,a_c\). Theorem~\ref{thm:main}, applied in
\(\overline S\) to every homogeneous ideal containing this sequence,
gives both the Eisenbud--Green--Harris and lex-plus-powers assertions for
\(\overline f_1,\ldots,\overline f_c\).

The Artinian reduction of Caviglia--Maclagan
\cite[Proposition~10]{CM08} transfers the Hilbert-function assertion to
\(f_1,\ldots,f_c\). Theorem~4.1 of Caviglia--Kummini \cite{CK14}
transfers the Betti assertion as well. Their formulation uses Betti
numbers of ideals. For any proper homogeneous ideal \(K\) in a standard
graded polynomial ring \(A\), the exact sequence
\(0\to K\to A\to A/K\to0\) gives
\[
 \beta^A_{p+1,j}(A/K)=\beta^A_{p,j}(K)
 \qquad(p\ge0,\ j\in\Z).
\]
The remaining quotient Betti number is \(\beta^A_{0,0}(A/K)=1\).
Thus their convention is equivalent to the quotient convention used
here. The lexicographic embedding modulo \(P_F\) in that reduction
selects exactly the prescribed \(q_d\) greatest monomials outside \(P_F\)
in each degree. A linear change of coordinates preserves Hilbert functions
and graded Betti numbers, so this proves the corollary over \(\C\).

Now let \(F\) be any characteristic-zero field. Choose homogeneous
generators \(g_1,\ldots,g_r\) of \(I\), and let \(F_0\subseteq F\) be
the field generated over \(\Q\) by all coefficients of the \(g_i\) and
the \(f_i\). Set
\[
 S_0=F_0[x_1,\ldots,x_n],\qquad
 I_0=(g_1,\ldots,g_r,f_1,\ldots,f_c)\subseteq S_0.
\]
Then \(I_0S_F=I\). The sequence \(f_1,\ldots,f_c\) is regular in
\(S_0\): after tensoring each multiplication map with the faithfully
flat extension \(F/F_0\), it is the corresponding injective map over
\(S_F\). The finitely generated characteristic-zero field \(F_0\)
admits an embedding into \(\C\). Extending by this embedding preserves
regularity, so the case just proved applies to
\[
 S_{\C}=\C\otimes_{F_0}S_0,\qquad I_{\C}=I_0S_{\C}.
\]
No embedding of \(F\) into \(\C\) is needed.

Let \(J_{\C}\) be the resulting lex-plus-powers ideal. It is a monomial
ideal, so let \(J_0\subseteq S_0\) be generated by the same monomials and
put \(J_F=J_0S_F\). For \(E=F\) or \(E=\C\), let
\(S_E=E\otimes_{F_0}S_0\). Graded field base change preserves the
dimensions of homogeneous pieces. It also gives, for \(M_0=S_0/I_0\)
and for \(M_0=S_0/J_0\),
\[
 E\otimes_{F_0}\Tor^{S_0}_p(M_0,F_0)_j
 \;\cong\;
 \Tor^{S_E}_p(E\otimes_{F_0}M_0,E)_j.
\]
For example, this follows by tensoring the Koszul complex of the
variables with \(E\); its graded homology commutes with this flat
extension. Consequently the Hilbert equalities and Betti inequalities
over \(\C\) give the same equalities and inequalities over \(F\).
Finally, the monomials of \(P_F\) and the lexicographic order are
determined by their exponents, while the dimensions defining \(q_d\)
are unchanged by these field extensions. Hence \(J_F\) is precisely
the degreewise ideal \(J\) prescribed in the corollary.
\end{proof}

The argument changes the coefficient field only through extensions of
the common characteristic-zero field \(F_0\). It makes no assertion for
arbitrary regular sequences in positive characteristic.
